\chapter*{Resumen}
\addcontentsline{toc}{chapter}{Resumen}
\markboth{Resumen}{Resumen}
\label{ch:resumen}

Esta tesis desarrolla un motor de estructura atómica, escrito desde cero en Julia, que resuelve las ecuaciones de Hartree-Fock restringido de capa abierta sobre una base de B-splines, y lo aplica a la secuencia homóloga $np^2$ del grupo 14: carbono, silicio, germanio y estaño. Las ecuaciones se discretizan con el método de Galerkin, los potenciales de interacción se obtienen de una ecuación de Poisson radial con una condición de Robin en el borde, y el ciclo autoconsistente se acelera con una extrapolación de Pulay sobre el conmutador proyectado, que en el germanio y el estaño reduce el número de iteraciones por factores de $2.8$ y $4.8$. Sobre los orbitales convergidos se añaden dos correcciones, ninguna con parámetros ajustables: la correlación de la pareja de valencia, mediante interacción de configuraciones con las capas internas congeladas, y el acoplamiento espín-órbita del Hamiltoniano de Breit-Pauli, diagonalizado de forma exacta en cada bloque de momento angular total.

El motor reproduce el límite de Hartree-Fock no relativista: las energías totales coinciden con las de Froese Fischer dentro de $3.3\times10^{-5}$ Hartree, las integrales de Slater dentro de $1.8\times10^{-5}$, y el cociente del virial se aparta de 2 en $1.5\times10^{-5}$ en el peor caso. El modelo completo reproduce el ordenamiento de los cinco niveles de la configuración en los cuatro elementos y sigue el tránsito del acoplamiento de Russell-Saunders al intermedio: la razón entre los dos intervalos del triplete fundamental cae de $1.91$ en el silicio a $1.14$ en el estaño, frente a $1.89$ y $1.03$ en el experimento. El peso del carácter $^1D_2$ en el nivel $^3P_2$ sigue la misma tendencia, aunque el modelo mezcla de menos en los elementos pesados: $5.5\%$ en el estaño, frente al $9.8\%$ que implica el factor $g$ medido.

Las discrepancias restantes se asocian a aproximaciones concretas del modelo. En el carbono, la razón de intervalos medida se aparta de la regla de Landé en un dieciocho por ciento, lo que apunta a los términos de dos cuerpos del Hamiltoniano de Breit-Pauli que el modelo omite. En el germanio y el estaño los intervalos del triplete quedan entre un 11 y un 19 por ciento por debajo de los medidos, y el parámetro de espín-órbita que exigen los niveles supera incluso la cota de núcleo desnudo: falta una contracción del orbital de valencia, que pueden producir tanto los efectos relativistas como la polarización de las capas internas, ambos fuera del modelo. El tratamiento consistente de los términos singlete, en particular de la casi degeneración $ns^2np^2 \leftrightarrow np^4$ que domina el error del $^1S_0$ del carbono, exige un cálculo de interacción de configuraciones de cuatro electrones.

\vspace{1em}
\noindent\textbf{Palabras clave:} Hartree-Fock restringido de capa abierta, B-splines, estructura fina, acoplamiento intermedio, interacción de configuraciones, secuencia $np^2$.

%%% Local Variables:
%%% mode: LaTeX
%%% TeX-master: "../main"
%%% End:
